\section{验证：从“能够运行”到可审计的等价性证据}
\label{sec:validation}

编译器验证最容易犯的认识论错误，是把一个必要条件误当成充分条件：语法正确不等于语义正确，能够链接不等于能够运行，若干实现给出相同答案也不等于答案正确。因而，本节不把“测试通过”压缩为一个布尔量，而把论证拆成四层：（1）由案例定义独立导出的答案；（2）每种表示自身的结构合法性；（3）汇编、链接和进程执行所要求的 ABI/ELF 不变量；（4）跨表示的可观察行为比较。这里的\emph{差分测试}（differential testing）指以多个独立实现的分歧作为缺陷信号；它善于发现错译，却不能在所有实现共同出错时给出警报。Csmith 的经验研究表明，在回避未定义及未指定行为之后，随机差分测试仍能在成熟 C 编译器中发现大量崩溃和静默错译~\cite{val-yang2011csmith}。这也解释了为什么本文既比较四种表示，又保留一个不依赖任何被测实现的算术预言机。

为避免把解释混入测量，以下采用四种证据标签。\observed 指命令输出或文件内容；\inferred 指由规范和观测共同推出、但没有被某个工具直接打印的结论；\textbf{假设：}指尚未由当前实验排除的解释；\textbf{建议：}指面向后续工作的设计选择。这一区分尤其重要：例如，目标文件中出现 \code{R\_RISCV\_CALL\_PLT} 是观测；链接器据此完成了对运行时调用的地址修补，则是结合最终符号表和 psABI 规则作出的推断~\cite{val-riscv-psabi}。

\subsection{答案先于程序：六输入预言机}

两个常量数组逐项相乘后得到
\[
 (-8,-3,7,12,-9,-10,12,15),
\]
再把每项限制到闭区间 $[-8,12]$，得到
\[
 t=(-8,-3,7,12,-8,-8,12,12).
\]
若输入为 $x$，有效长度为 $n=\min(8,\max(0,x))$，则规范结果就是前缀和
\begin{equation}
  F(x)=\sum_{i=0}^{n-1}t_i .
  \label{eq:validation-oracle}
\end{equation}
这个推导只使用案例中的数学契约，不读取任何可执行文件的输出。因此，表~\ref{tab:validation-oracle} 中的字节串是\emph{已知答案测试}（known-answer test）的预言机，而非先运行 C 程序再抄录的“黄金输出”。

\begin{table}[t]
  \centering
  \caption{独立推导的六输入预言机；\code{\textbackslash n} 表示单个 LF 字节。}
  \label{tab:validation-oracle}
  \small
  \begin{tabularx}{\textwidth}{@{}r r l X@{}}
    \toprule
    输入 $x$ & $n$ & 标准输出 & 覆盖意图 \\
    \midrule
    $-1$ & 0 & \code{0\textbackslash n} & 下界截断成立；循环零次 \\
    0    & 0 & \code{0\textbackslash n} & 精确下边界；循环零次 \\
    1    & 1 & \code{-8\textbackslash n} & 一次迭代、负数乘法与下界 \\
    4    & 4 & \code{8\textbackslash n} & 内部长度、多次迭代和累加 \\
    8    & 8 & \code{16\textbackslash n} & 精确上边界、访问全部元素 \\
    10   & 8 & \code{16\textbackslash n} & 上界截断成立；不发生越界 \\
    \bottomrule
  \end{tabularx}
\end{table}

每次执行还必须满足返回码为 0。仓库提供的 \code{sylib.c} 通过析构函数在标准错误输出固定写入 \code{TOTAL: 0H-0M-0S-0us\textbackslash n}；这是运行时的计时诊断，不是案例输出。故测试器应分别比较三个通道：逐字节比较表中的标准输出，精确接受该标准错误尾注，并单独检查退出状态，而不能用删除所有空白、合并两个流或仅比较屏幕文本的办法“放宽”预言机。

\observed 在 Ubuntu/WSL 环境中，C 观察载体、以 C 前端承载的 SysY 源、手写 LLVM IR 和手写 RV64 汇编共四个静态 RV64 可执行文件，对表中六个输入进行了 $4\times6=24$ 次 QEMU 执行；24 次的标准输出均逐字节匹配，退出状态均为 0，标准错误均为上述固定尾注。由此可说“这 24 个样本未发现行为差异”，不能说“这四个程序对所有输入均等价”。这种措辞差异不是谦辞，而是有限实验与全称命题之间的逻辑边界。

\subsection{分层核验四种表示}

\paragraph{C 观察载体与 SysY。}
纯 SysY 不包含 C 预处理器；因此 \code{clipped\_dot.c} 的角色不是冒充另一份语言规范，而是让预处理、Clang AST、自动生成 IR 和后端产物可以被观测的\emph{载体}。\observed 在本次便携捕获中，\code{VECTOR\_LENGTH} 与 \code{TERM\_MIN} 等宏名已从 \code{clipped\_dot.i} 消失，数组界变为 8；JSON AST 的根为 \code{TranslationUnitDecl}，并包含 \code{clipped\_dot}、\code{WhileStmt}、\code{IfStmt} 和 \code{CallExpr}。同一脚本用固定目标 \code{riscv64-unknown-linux-gnu/rv64gc/lp64d} 生成 $O0$ 与 $O2$ 的 IR 和汇编，并确认两份 IR 不逐字节相同。这里的强结论仅是“阶段产物存在且含预期结构”，不是“优化必然正确”。LLVM 自身也把小型回归测试、单元测试和整程序参考输出视为不同层次的证据~\cite{val-llvm-testing}。

\paragraph{手写 LLVM IR。}
该模块明确给出 RV64 数据布局和目标三元组；两个 \code{[8 x i32]} 常量保存数组初值；\code{clipped\_dot} 用两个 $\phi$ 节点表示循环携带的索引和累加值，以带符号比较、\code{select} 和 $i32$ 的 \code{mul}/\code{add} 表示逐项截断。\code{main} 用 \code{getelementptr} 的等价内存布局、\code{llvm.memcpy}、长度夹取和三项运行时声明连接到外界。\observed LLVM 18.1.3 的 \code{llvm-as} 将其汇编为位码，\code{opt -passes=verify} 接受该位码，Clang 又将原始 \code{.ll} 编译为 2480 字节的 RV64 可重定位目标。验证器证明的是 IR 结构满足 LLVM 的良构约束；随后六组执行才检查所选路径的结果，两者不可互相替代。LLVM 的具体语义以 LangRef 为准，例如 \code{inbounds} 和毒值传播不能仅凭“看起来像地址计算”理解~\cite{val-llvm-langref}。

\paragraph{手写 RV64 汇编。}
核心函数是叶函数，以 \code{lw} 取有符号 32 位元素，以 \code{mulw}/\code{addw}/\code{addiw} 保持 SysY \code{int} 的低 32 位运算，并用带符号分支实现上下界；\code{main} 分配 80 字节栈帧、在偏移 72 保存返回地址，在所有调用点维持 16 字节栈对齐，返回前恢复二者。\observed GNU 汇编器接受此文件并产生 2152 字节目标文件；反汇编仍可定位循环回边、四字节缩放寻址与运行时调用。\inferred 由于核心函数只使用调用者保存的 \code{t0--t5}/\code{a0--a4} 且不再调用其他函数，无须额外保存被调用者保存寄存器；这依赖 RISC-V ELF psABI 的寄存器约定，而不是由一次成功执行证明~\cite{val-riscv-psabi}。

表~\ref{tab:validation-artifacts} 仅报告本次捕获目录中实际文件的字节数。它们适合发现产物缺失或意外剧变，不适合推出性能：静态链接可执行文件的大部分空间来自启动代码、libc 和运行时，$O2$ 汇编更短也不自动意味着在真实硬件上更快。

\begin{table}[t]
  \centering
  \caption{本次捕获中直接观测到的文件尺寸（字节）。}
  \label{tab:validation-artifacts}
  \small
  \begin{tabular}{@{}l r r@{}}
    \toprule
    产物 & 字节数 & 解释边界 \\
    \midrule
    C 载体 $O0$ / $O2$ LLVM IR & 8241 / 7417 & 文本尺寸，不是动态指令数 \\
    C 载体 $O0$ / $O2$ RV64 汇编 & 5044 / 3128 & 生成形式，不含最终链接代码 \\
    手写 IR / 手写汇编目标 & 2480 / 2152 & 均为 \code{ET\_REL} \\
    C / SysY 静态可执行文件 & 626832 / 626832 & 含 glibc 与 SysY 运行时 \\
    手写 IR / 汇编静态可执行文件 & 626776 / 626664 & 同上；尺寸差异不作性能结论 \\
    \bottomrule
  \end{tabular}
\end{table}

\subsection{跨越汇编器、链接器与进程边界}

只运行最终程序会掩盖“错误的库恰好没有被调用”等问题，因此验证必须保留链接前后的证据链。\observed 四个输入目标的 ELF 头均报告 \code{ELF64}、小端、\code{ET\_REL} 和机器 \code{RISC-V}。以手写汇编目标为例，\code{getint}、\code{putint}、\code{putch} 是未定义全局符号；调用点具有 \code{R\_RISCV\_CALL\_PLT} 及相邻 \code{R\_RISCV\_RELAX} 重定位。最终文件则为 \code{ET\_EXEC}，具有程序头和非零入口点（手写汇编版本为 \code{0x10664}），运行时符号已进入静态映像。于是，能够陈述的完整链条是：源文件留下未解析调用，汇编器将其编码为符号加重定位，链接器从归档抽取成员并完成解析，QEMU 才执行所得 Linux/RV64 进程映像。

运行时还暴露了一个容易被“同为 RISC-V”掩盖的接口边界。\observed 随附 \code{libsysy\_riscv.a} 的未定义符号包含 \code{\_impure\_ptr}，而 Debian/Ubuntu 的 glibc RV64 链路并不提供该 newlib 状态符号。实验没有伪造一个占位定义，而是用同一 \code{rv64gc/lp64d} 工具链重新编译仓库的 \code{sylib.c}，产生 19444 字节的 \code{libsysy-glibc-rv64.a}，再由交叉 GCC 驱动器以 \code{-static -no-pie} 链接。重编译还如实产生了若干 \code{-Wunused-result} 警告，因为该教学运行时没有检查 \code{scanf} 返回值；六组合法整数输入没有触发输入失败，但这不是可忽略一般输入错误的证明。\inferred 这说明兼容性至少包含 ISA、ELF ABI、C 库环境和运行时符号契约四层；仅查看机器类型相同不足以判定两个目标能够安全链接。

\subsection{证据强度、可复现性与尚未证明之处}

本实验采用“独立答案 $+$ 多实现差分 $+$ 结构不变量”的三角测量，而非把同一工具生成的两个文件当作相互独立的见证。方法的强项是故障定位：若 IR 不能通过验证器，问题位于 IR 良构性；若目标缺少重定位，问题在汇编边界；若只有一个可执行文件输出不同，差分矩阵能定位可疑表示。它的局限同样明确。

第一，六个输入是刻意选择的边界覆盖，不是形式证明；数组内容固定，且没有穷举所有 $2^{32}$ 个输入。第二，C/SysY 是由前端生成的形式，LLVM IR 与 RV64 汇编是人工编写的形式；四者的“独立性”并不对称，因为它们最后仍共享链接器、运行时和 QEMU。第三，\code{opt} 的验证器检查结构约束，不证明手写 IR 与 SysY 源语义等价；QEMU 执行也不验证真实 RV64 微体系结构的时序或性能。第四，本节的数字绑定于表~\ref{tab:validation-env} 的环境；更换 LLVM 主版本、glibc、binutils 或链接选项可能改变文本、尺寸、重定位和地址，而不改变程序语义。

\begin{table}[t]
  \centering
  \caption{产生本节完整 RV64 观测的工具环境；版本是实验输入，而非装饰性元数据。}
  \label{tab:validation-env}
  \small
  \begin{tabularx}{\textwidth}{@{}l X@{}}
    \toprule
    组件 & 观测版本/配置 \\
    \midrule
    主机边界 & Linux 6.18.33.2 WSL2，x86-64；Python 3.12.3 \\
    LLVM & Ubuntu Clang/LLVM/\code{llvm-as}/\code{opt} 18.1.3 \\
    GNU 交叉工具链 & \code{riscv64-linux-gnu-gcc} 13.3.0；binutils 2.42 \\
    执行器 & \code{qemu-riscv64} 8.2.2 \\
    目标/链接 & \code{riscv64-unknown-linux-gnu}，\code{rv64gc/lp64d}，静态非 PIE \\
    \bottomrule
  \end{tabularx}
\end{table}

\paragraph{与编译器验证研究的关系。}
\emph{翻译验证}（translation validation）不预先证明整个编译器正确，而是对每一次具体翻译检查源程序与目标程序的对应关系~\cite{val-pnueli1998translation}。Alive2 把这一思想用于 LLVM IR，以 SMT 求解器检查优化前后程序的精化关系；其部署曾发现 LLVM 中的新缺陷，但循环展开等资源界限意味着它仍可能漏报~\cite{val-lopes2021alive2}。因此，\textbf{建议：}后续可把案例的 $O0/O2$ IR 对交给 Alive2，作为六输入执行之外的有界翻译验证；若再加入无未定义行为的固定种子生成器，则可扩大测试分布。两者都应是新增证据层，而非把“bounded”误写为全程序、全输入的等价证明。更广泛的形式验证与编译器前沿留待第~12 节讨论。

\begin{table}[t]
  \centering
  \caption{从干净检出复现实验的检查清单。}
  \label{tab:validation-checklist}
  \small
  \begin{tabularx}{\textwidth}{@{}p{0.05\textwidth}X@{}}
    \toprule
    & 可复现性检查项 \\
    \midrule
    $\square$ & 记录提交哈希、工作树状态、操作系统、Python、Clang/LLVM、交叉 GCC/binutils 与 QEMU 版本。\\
    $\square$ & 从 \code{preflight/src}、\code{preflight/ir}、\code{preflight/asm} 重新捕获产物；不复用未知来源的目标文件。\\
    $\square$ & 对手写 IR 运行 \code{llvm-as} 与 \code{opt -passes=verify}；对手写汇编实际汇编，而非只做文本匹配。\\
    $\square$ & 检查各目标为 ELF64、小端、RISC-V、\code{ET\_REL}，并保留符号、重定位和反汇编。\\
    $\square$ & 用目标工具链重建运行时；记录随附归档的 \code{\_impure\_ptr} 不兼容证据，不以宿主库替代。\\
    $\square$ & 检查最终文件为带入口和程序头的 \code{ET\_EXEC}；确认运行时符号从未定义变为已解析。\\
    $\square$ & 对四种形式执行六输入矩阵；分别逐字节比较 stdout、stderr 和退出状态。\\
    $\square$ & 保存 \code{capture.json} 中的参数边界与命令序列；报告任何警告、跳过项和工具版本差异。\\
    \bottomrule
  \end{tabularx}
\end{table}

最终，本文的经验结论是有限而明确的：在记录的工具链上，四条实现路径都跨过了其各自的语法/结构、ABI、重定位、静态链接和执行边界，并在六个独立给定答案上得到一致的可观察结果。可证伪条件也同样明确：任一输入产生不同字节、非零退出状态、非预期诊断，或任一结构不变量消失，都会推翻相应层次的结论。这样的证据链比一句“编译成功”更繁琐，却正是把编译系统从黑箱变成可检查对象所需的代价。
